\begin{table}[!htbp]
  \centering
  \small
  \setlength{\tabcolsep}{4pt}
  \begin{threeparttable}
  \begin{tabular}{l l l}
    \toprule
    \textbf{维度} & \textbf{KIVI-style}~\cite{liu2024kivi} & \textbf{INT4-RoleAlign}（本文） \\
    \midrule
    \textbf{Cache format} &
      per-channel K + per-token V &
      \textbf{同上}（完全 format-compatible） \\
    \addlinespace[2pt]
    \textbf{Calibration} &
      runtime absmax/min &
      \textbf{离线 K/V 分路径校准}，behavior 目标 \\
    \addlinespace[2pt]
    \textbf{Scale 固化} &
      每次 prefill/decode 重算 &
      \textbf{校准产物 JSON}，推理时只读 \\
    \addlinespace[2pt]
    \textbf{Allocator} &
      无（所有层统一 bit 分配） &
      可选扩展至 behavior-guided allocator（§\ref{sec:exp-cross-model}） \\
    \addlinespace[2pt]
    \textbf{产物接口} &
      运行时统计范围更新，无离线校准产物 &
      \textbf{可固化 / 可审计}，适合合规部署 \\
    \bottomrule
  \end{tabular}
  \begin{tablenotes}
    \footnotesize
    \item 同 format 不同 calibration：两者的 KV cache 布局完全一致，解码语义相同，
    \textbf{唯一差异}在于量化 scale 的来源。
    \item 定位：本文主张 INT4-RoleAlign 是 behavior-guided framework 在 KIVI 非对称 format 上的
    落地实例；两者之间是 \emph{calibration philosophy}
    的对比，而非 \emph{format / 架构} 的竞争。
  \end{tablenotes}
  \end{threeparttable}
  \caption{INT4-RoleAlign 与 KIVI-style 设计差异。共享量化格式、差异集中在 calibration；
  其中 RoleAlign 的 K-path 与 V-path 分别由注意力分布代理和输出扰动代理校准；Allocator 扩展与可审计产物接口是本文在 KIVI 基础上新增的两个维度。}
  \label{tab:app-rolealign-vs-kivi}
\end{table}
